% ==============================================================================
% Appendix
% ==============================================================================

\chapter*{Appendix}
\addcontentsline{toc}{chapter}{Appendix}
\markboth{Appendix}{Appendix} 

The following appendices provide supplementary technical documentation for the Smart Agriculture System, including the user manual for the mobile application, hardware setup instructions, dataset parameters, AI model details, and troubleshooting guidelines.

% ==============================
% Appendix A
% ==============================

\section*{Appendix A: User Manual and Troubleshooting}
\addcontentsline{toc}{section}{Appendix A: User Manual and Troubleshooting}

\subsection*{A.1 Mobile Application User Manual}
The Smart Agriculture mobile application is designed for ease of use and real-time monitoring. 
\textbf{Account Registration:} Open the application and navigate to the `Register' screen. Provide your full name, email address, phone number, and a secure password.

\textbf{Field Management:} Upon logging in, the Dashboard displays your registered agricultural fields. To add a new field, tap the `+' floating action button and enter the field details (such as crop type and area).

\textbf{Sensor Monitoring:} Tap on a specific field to view real-time telemetry. The detailed dashboard displays a live moisture gauge chart, temperature and humidity readouts, and the latest update timestamp.

\textbf{Irrigation Control:} Navigate to the `Irrigation' tab to toggle the water pump manually or set up automated scheduling based on predetermined soil moisture thresholds.

\textbf{AI Recommendations:} Navigate to the `Recommendations' tab to view crop suggestions generated by the AI model based on your field's historical environmental data.



\subsection*{A.2 Troubleshooting Guide}
\textbf{ESP32 Offline Status:} Check the Wi-Fi connection in the field. Ensure the power supply (battery or adapter) is properly connected and providing adequate voltage. Restart the microcontroller if necessary.

\textbf{Inaccurate Sensor Readings:} For the DHT22 sensor, ensure it is not exposed to direct sunlight or rain. For the soil moisture sensor, ensure the probes are fully inserted into the soil and periodically clean off any corrosive buildup on the metal contacts.

\textbf{Pump Not Activating:} Check the relay module indicator light. Verify that the 12V power supply for the pump is active and that the GPIO signal wires are securely connected to the ESP32.



% ==============================
% Appendix B
% ==============================
\section*{Appendix B: Hardware Configuration and AI Model Details}
\addcontentsline{toc}{section}{Appendix B: Hardware Configuration and AI Model Details}

\subsection*{B.1 Hardware Setup and Pin Configuration}
The Field Unit is centered around an ESP32 microcontroller. The sensors and actuators are connected to the following GPIO pins to ensure reliable data acquisition and control:
\begin{table}[H]
\centering
\begin{tabular}{|l|l|l|}
\hline
\textbf{Component} & \textbf{Pin Type} & \textbf{ESP32 GPIO Pin} \\ \hline
Soil Moisture Sensor & Analog Input & GPIO 34 \\ \hline
DHT22 (Temperature/Humidity) & Digital Input & GPIO 4 \\ \hline
LDR (Light Sensor) & Analog Input & GPIO 35 \\ \hline
Rain Sensor & Digital Input & GPIO 15 \\ \hline
Water Flow Sensor & Digital Input & GPIO 2 \\ \hline
Relay Module (Pump Control) & Digital Output & GPIO 5 \\ \hline
\end{tabular}
\caption{ESP32 Pin Configuration for Field Unit}
\label{tab:pin_config}
\end{table}
Ensure all common ground connections are shared between the ESP32, the sensor modules, and the relay module. The relay typically requires a separate 5V logic supply to safely isolate the pump's higher voltage circuit from the microcontroller.

\subsection*{B.2 Dataset and AI Model Parameters}
The crop recommendation engine is powered by a machine learning model developed in Python using the Scikit-Learn library. The core details are as follows:
\textbf{Dataset Features:} Soil Moisture (\%), Temperature ($^\circ$C), Humidity (\%), Soil Type (encoded via LabelEncoder), and Season (encoded via LabelEncoder).

\textbf{Model Type:} Random Forest Classifier.

\textbf{Hyperparameters:} \texttt{n\_estimators} = 150, \texttt{max\_depth} = 10, \texttt{min\_samples\_leaf} = 3.

\textbf{Validation and Confidence Score:} The model achieved robust accuracy during 5-fold cross-validation on an 80/20 train-test split. The confidence score returned by the API is calculated mathematically by determining the percentage of the 150 decision trees that voted for the majority crop class ($C = \max_c P(c|x) \times 100\%$).

